\begin{lemma}\label{lem:odd-root-three-disjoint}
If an oriented graph $R$ on three or four vertices has positive
outdegree at every vertex, it contains three vertices whose
outneighborhoods in $R$ are pairwise disjoint.
\end{lemma}
\begin{proof}
Complete $R$ to a tournament; this preserves positive outdegree,
and any disjointness proved in the completion holds in $R$.
On three vertices the completion must be a directed triangle.
On four vertices, if some vertex has outdegree three, the remaining
three vertices form a directed triangle and their outneighborhoods
are disjoint.  Otherwise the tournament has score multiset
$(1,1,2,2)$.  Let $a,b$ have outdegree one and orient their mutual
arc as $a\to b$.  If the other vertices are $c,d$, then both point
to $a$, while $b$ points to exactly one of them, say $c$.  Thus
$d\to b$.  The arc $d\to c$ would give $d$ outdegree three, so
$c\to d$.  The outneighborhoods of $a,b,c$ are respectively
$\{b\},\{c\},\{a,d\}$, which are pairwise disjoint.
\end{proof}
